The CMS experiment at the LHC is upgrading its endcap
calorimeters with the High-Granularity Calorimeter (HGCAL). Long-lived particles can decay far from the
interaction point and produce photons whose trajectories do not point back to
it. Reconstructing these photons is a challenge for \cluethreed{}, the
clustering algorithm developed for HGCAL.

This project studies \cluethreed{} performance using simulated single- and
two-photon events, with and without pileup. We improved the displaced-photon
generator and added displacement observables to the standard validation tools.
The results show that \cluethreed{} loses efficiency as displacement increases,
with single showers fragmenting and nearby showers becoming mixed. Adding pileup
leaves the signal-photon efficiency largely unchanged. Tests of the alternative
\cluestering{} algorithm show that suitable configurations improve
displaced-photon reconstruction while keeping nearby showers separate. Further
tuning and tests with realistic long-lived-particle decays are needed.
